\documentclass[12pt,a4paper]{article}

\usepackage[brazil]{babel}
\usepackage{fontspec}
\usepackage{xeCJK}
\usepackage{xcolor}
\usepackage{ruby}
\usepackage{amsmath}
\usepackage{tcolorbox}
\usepackage{enumitem}
\usepackage{tabularx}
\usepackage{booktabs}

\setmainfont{Latin Modern Roman}
\setCJKmainfont{Noto Serif CJK JP}

\definecolor{jpnavy}{RGB}{20, 35, 60}
\definecolor{jpred}{RGB}{180, 40, 40}
\definecolor{lighthight}{RGB}{245, 247, 250}

\pagestyle{empty}

\begin{document}

\begin{center}
    {\LARGE\bfseries 第十一課}\\[4pt]
    \small Clima, Estações e Conversa Casual
\end{center}

\vspace{0.3cm}
\hrule height 1pt
\vspace{0.5cm}

\begin{tcolorbox}[colback=lighthight, colframe=jpnavy, title=\textbf{Dica Cultural \& Estações do Ano}, fonttitle=\bfseries]
\small
As quatro estações (\textit{Shiki}) são centrais na cultura japonesa, influenciando a culinária, festivais e hábitos do dia a dia. Comentários sobre o tempo (\textit{Aisatsu}) são a forma mais comum e sutil de iniciar uma conversa sociável.
\end{tcolorbox}

\vspace{0.2cm}

\subsection*{Estações, Clima e Temperaturas (\textit{Tenki to Kisotsu})}

\vspace{0.2cm}

\noindent
\small
\begin{tabularx}{\textwidth}{>{\bfseries\color{jpnavy}}l X X}
\toprule
Categoria & Kanji / Kana (Romaji) & Tradução em Português \\
\midrule
Estações & 春 (\textit{Haru}) / 夏 (\textit{Natsu}) & Primavera / Verão \\
& 秋 (\textit{Aki}) / 冬 (\textit{Fuyu}) & Outono / Inverno \\
\midrule
Condições & 天気 (\textit{Tenki}) / 晴れ (\textit{Hare}) & Clima, Tempo / Tempo limpo \\
& 雨 (\textit{Ame}) / 雪 (\textit{Yuki}) / 曇り (\textit{Kumori}) & Chuva / Neve / Nublado \\
\bottomrule
\end{tabularx}

\vspace{0.3cm}

\noindent
\small
\begin{tabularx}{\textwidth}{>{\bfseries\color{jpnavy}}l X X}
\toprule
\multicolumn{3}{c}{\textbf{\color{jpnavy}Temperaturas: Clima vs. Coisas / Alimentos}} \\
\midrule
Clima Quente & 暑い (\textit{Atsui}) & 今日は暑いですね。 (Hoje está quente, né?) \\
Clima Frio & 寒いです (\textit{Samui}) & 今日は寒いですね。 (Hoje está frio, né?) \\
\midrule
Coisa / Comida Quente & 熱い (\textit{Atsui}) & 熱いお茶 (\textit{Atsui o-cha} - Chá quente) \\
Coisa / Comida Fria & 冷たい (\textit{Tsumetai}) & 冷たい水 (\textit{Tsumetai mizu} - Água gelada) \\
\bottomrule
\end{tabularx}

\newpage

\begin{center}
    {\LARGE\bfseries 会話 (Diálogo: Reclamando do Clima)}
\end{center}

\textit{Cenário: Kenji e Yuta caminham pelo campus em um dia ensolarado no meio do verão.}

\vspace{0.8em}

\textbf{\ruby{健二}{けんじ}:} うわあ、今日は本当に暑いですね！\\
\textit{Kenji:} Uwaa, kyou wa hontou ni atsui desu ne!\\
Kenji: Uau, hoje está quente de verdade, né?

\vspace{0.4em}

\textbf{\ruby{裕太}{ゆうた}:} 本当ですね。日本の夏は蒸し暑いですから、大変ですよ。\\
\textit{Yūta:} Hontou desu ne. Nihon no natsu wa mushiatsui desu kara, taihen desu yo.\\
Yuta: É mesmo! O verão no Japão é muito abafado, então é complicado.

\vspace{0.4em}

\textbf{\ruby{健二}{けんじ}:} ええ、汗をたくさんかきました。昨日は雨でしたね。\\
\textit{Kenji:} Ee, ase wo takusan kakimashita. Kinou wa ame deshita ne.\\
Kenji: É, suei bastante. Ontem choveu, né?

\vspace{0.4em}

\textbf{\ruby{裕太}{ゆうた}:} そうですね。昨日は少し涼しかったですが、今日は全然違いますね。\\
\textit{Yūta:} Sou desu ne. Kinou wa sukoshi suzushikatta desu ga, kyou wa zenzen chigaimasu ne.\\
Yuta: Pois é. Ontem estava um pouco fresco, mas hoje está totalmente diferente, né?

\vspace{0.4em}

\textbf{\ruby{健二}{けんじ}:} 何か冷たいものを飲みませんか？ジュースとか...\\
\textit{Kenji:} Nani ka tsumetai mono wo nomimasen ka? Jūsu toka...\\
Kenji: Que tal bebermos algo gelado? Um suco ou algo assim...

\vspace{0.4em}

\textbf{\ruby{裕太}{ゆうた}:} いいですね！あそこの自動販売機で買いましょう！\\
\textit{Yūta:} Ii desu ne! Asoko no jidouhanbaiki de kaimashou!\\
Yuta: Ótima ideia! Vamos comprar ali na máquina de vendas!

\newpage

\subsection*{Vocabulário Útil: Clima e Sensações (\textit{Tango})}

\vspace{0.2cm}

\noindent
\small
\begin{tabularx}{\textwidth}{>{\bfseries\color{jpnavy}}l X X}
\toprule
Palavra (Romaji) & Kanji / Kana & Tradução em Português \\
\midrule
Mushiatsui & 蒸し暑い (むしあつい) & Abafado / Úmido \\
Suzushii & 涼しい (すずしい) & Agradável / Fresco (clima) \\
Warm / Natakai & 暖かい (あたたかい) & Agradável / Morno (clima/comida) \\
Kaze & 風 (かぜ) & Vento \\
Kaze ga fuku & 風が吹く (かぜがふく) & Ventar \\
Ame ga furu & 雨が降る (あめがふる) & Chover \\
\midrule
Jidouhanbaiki & 自動販売機 (じどうはんばいき) & Máquina de venda automática \\
Tsumetai mono & 冷たいもの & Coisa / Bebida gelada \\
Atsui mono & 熱いもの & Coisa / Bebida quente \\
Ase & 汗 (あせ) & Suor \\
\bottomrule
\end{tabularx}

\vspace{0.5cm}

\begin{tcolorbox}[colback=jpred!5, colframe=jpred, title=\textbf{Dinâmica em Grupo}, fonttitle=\bfseries]
\small
Pratique conversas de "elevador/corredor" com os colegas usando as expressões de clima e a partícula confirmatória \textbf{ne} ("né?").

\vspace{0.2cm}

\begin{enumerate}[leftmargin=*]
    \item \textbf{Puxar conversa sobre o tempo:} \textit{Kyou wa [Atsui / Samui / Mushiatsui] desu ne!} (Hoje está [Quente / Frio / Abafado], né?)
    \item \textbf{Concordar com o colega:} \textit{Sou desu ne! Hontou ni [Adjetivo] desu ne.} (Pois é! Realmente está [Adjetivo], né.)
    \item \textbf{Sugerir uma pausa / bebida:} \textit{Tsumetai mono / Atsui mono wo nomimashou ka?} (Vamos tomar algo gelado / quente?)
\end{enumerate}
\end{tcolorbox}

\vspace{0.4cm}

\begin{center}
    \small\color{gray} \textit{お疲れ様でした！}
\end{center}

\end{document}